\chapter*{คำนำ}
\addcontentsline{toc}{chapter}{คำนำ}

คู่มือเล่มนี้เริ่มต้นจากคำถามที่เรียบง่ายแต่มีนัยสำคัญว่า \enquote{เราจะช่วยให้นักวิจัยคนหนึ่งมองเห็นคุณค่าของงานตนเอง และพาคุณค่านั้นไปถึงผู้คนได้อย่างไร} คำตอบมิได้อยู่ที่คำแนะนำเพียงครั้งเดียว แต่อยู่ที่ความสัมพันธ์ซึ่งทำให้ Mentee กล้าคิดชัดขึ้น เห็นหลักฐานตรงขึ้น และตัดสินใจด้วยความรับผิดชอบมากขึ้น

เนื้อหาเดิมจากการอบรมได้รับการเรียบเรียงใหม่เป็น handbook สำหรับการใช้งานระยะยาว มิใช่เอกสารประกอบสไลด์ แต่เป็นทั้งกรอบคิด คู่มือสนทนา และสมุดฝึกปฏิบัติ ท่านสามารถอ่านตามลำดับเพื่อสร้างพื้นฐาน หรือเปิดใช้เฉพาะบทก่อนเข้าสู่การสนทนากับ Mentee

\section*{วิธีใช้คู่มือ}
Part I ช่วยให้ท่านนิยามตัวตนและขอบเขตของบทบาท Mentor; Part II ใช้ Tech Portfolio เป็นภาษากลางในการมองเห็นการเติบโต; Part III แสดงวิธีใช้ AI อย่างรับผิดชอบ; Part IV พัฒนาทักษะมนุษย์ที่ทำให้บทสนทนามีคุณภาพ; Part V วางจังหวะการทำงานสามครั้งแรก; และ Part VI เชื่อมการ Mentoring รายบุคคลเข้ากับระบบนิเวศมหาวิทยาลัยแม่โจ้

\begin{mentornote}
คู่มือนี้เสนอกรอบ มิได้เสนอสูตรสำเร็จ การใช้เครื่องมือทุกชิ้นควรยึดเป้าหมายของ Mentee บริบทของงาน และจริยธรรมของความสัมพันธ์เป็นหลัก
\end{mentornote}

\section*{คำศัพท์ที่ใช้}
หนังสือใช้คำว่า \textbf{Mentor} สำหรับผู้ร่วมเดินทางที่ใช้ประสบการณ์ คำถาม และเครือข่ายเพื่อสนับสนุนการเติบโต และใช้คำว่า \textbf{Mentee} สำหรับเจ้าของเป้าหมายและเจ้าของการตัดสินใจ คำภาษาอังกฤษที่จำเป็นคงไว้เพื่อให้สอดคล้องกับการใช้งานในระบบนิเวศนวัตกรรม
